\section{Industrial Internet of Things (IIoT) in Collaborative Robotics}

\subsection{Evolution of Industrial Automation}
The evolution of industrial robotics in 2026 has transitioned from rigid, pre-programmed automation to adaptive, operational intelligence driven by the Industrial Internet of Things (IIoT) (ESA Automation, 2026). Within the context of collaborative robotic systems, IIoT serves as the foundational architecture for real-time data synchronization between physical manipulators and their digital counterparts. Modern frameworks increasingly prioritize Edge Computing to facilitate ultra-low latency decision-making, which is critical for safety in human-robot interaction where cloud-based delays are unacceptable (viAct, 2026). Furthermore, the integration of standardized communication protocols, such as OPC UA over Time-Sensitive Networking (TSN), has enabled a deterministic 'nervous system' for factories, allowing complex motion control data to share physical infrastructure with high-level telemetry without interference (Mieruński, 2026). These advancements support the deployment of Digital Twins, which utilize continuous sensor streams to enable predictive maintenance and autonomous process optimization (Xu et al., 2023; Lopes, 2025).


\subsection{The Role of IIoT in Collaborative Robotics and Industry 4.0}

The integration of the Industrial Internet of Things (IIoT) marks the transition from Industry's focus on mass automation to the Industry 4.0 paradigm, which emphasizes meaningful human-robot collaboration (HRC). In collaborative robotics, IIoT serves as more than a data conduit; it acts as a 'shared sensory awareness' layer. By leveraging high-speed communication and Edge Computing, IIoT-enabled systems can process environmental data from vision sensors and tactile feedback in real-time. This allows collaborative manipulators to adapt their trajectories and velocities dynamically based on human proximity, ensuring safety without the need for physical barriers. Furthermore, the ability to aggregate data from multiple distributed nodes enables a collective intelligence, where learning from one assembly task can be propagated across the entire network, significantly reducing the downtime associated with reprogramming traditional industrial arms.